\documentclass[10.5pt,a4paper]{article}
\usepackage[margin=22mm]{geometry}
\usepackage{fontspec}
\setmainfont{Arial}
\usepackage[hidelinks]{hyperref}
\usepackage{parskip}
\setlength{\parskip}{0.75em}
\pagestyle{empty}
\begin{document}

\begin{flushright}
Luciano Moffatt\\
Instituto de Química Física de los Materiales, Medio Ambiente y Energía (INQUIMAE), CONICET\\
Facultad de Ciencias Exactas y Naturales, Universidad de Buenos Aires\\
Ciudad de Buenos Aires, C1428EHA, Argentina\\
\href{mailto:lmoffatt@qi.fcen.uba.ar}{lmoffatt@qi.fcen.uba.ar}\\[0.6em]
30 August 2026
\end{flushright}

Dear Editors,

I am submitting ``Likelihood approximations distort the ion channel kinetic information in
macroscopic currents'' for consideration as a Research Article.

Kinetic mechanisms of ion channels are beginning to be chosen by comparing the Bayesian evidence
for rival schemes fitted to macroscopic currents, as we did for nine P2X$_2$ schemes in 2025. How
much the likelihoods behind such a comparison misreport the information a recording holds, and in
which directions, has never been measured for macroscopic currents. This manuscript measures both
for a ladder of eight likelihoods, from least squares to a filter that uses the channel state at
both ends of each sampling interval, against ten thousand exact simulations at every combination
of channel number, noise and sampling interval tested.
Every likelihood that ignores the correlation between successive samples misreports its
information by an order of magnitude, two- to threefold on the error bar, wherever that correlation
is present in the record, and no single rescaling repairs it, because the distortion differs
between parameter directions. The filter conditioning on both ends reports the uncertainty it
delivers, which I call calibrated, except where single openings become resolvable in the record,
and there its departure follows one measured power law. No measured condition gives both a
calibrated least-squares error bar and information about the single-channel current; the two lie
about a factor of a hundred apart in noise, and still a factor of three apart under the paper's
looser criterion.

In our 2025 ranking of nine P2X$_2$ schemes the winning mechanism changed when the likelihood's
recursion was removed, overturning a factor of six in evidence, and we published that
sensitivity. In the two-state scheme this manuscript shows that the likelihood without the
recursion, our control, reports at one representative condition about twenty times the
information the record holds, while the likelihood we argued for is calibrated there; whether
that carries to nine schemes is suggested and not shown, and the evidence is not recomputed here.

The measurement needs no fit to a real recording. At known parameters a correct likelihood
satisfies two identities: its score, the gradient of the log-likelihood, averages to zero, and the
covariance of the score equals the Fisher information the likelihood itself reports. The
calibrated filter was tested over 560 such combinations, all eight likelihoods over 210 of them.
The scheme has two states at an open probability of one half, the smallest that isolates the
approximation's own error; the failure comes from a correlation the likelihood does not carry,
which richer schemes share but are not tested here.

The closest prior check on a filter of this kind appeared in eLife, by M\"unch and colleagues in
2022. They counted how often the true value fell inside the reported interval over repeated fits,
a check that costs a fit per simulation and so covers a line of conditions. The two identities
cost one likelihood evaluation per simulation and return a matrix over a plane: which parameter
is misreported, and whether at each sample or across the record. The filter itself was published
with the P2X$_2$ analysis and has the form of a 1988 integrated-measurement Kalman filter; what is
new is its score and Fisher information, and the measurement built on them.

The engine with its exact simulator, the configurations and
the notebooks are at 10.5281/zenodo.22168409, the simulation output at 10.5281/zenodo.22167744,
and the likelihood with its score, Fisher information and diagnostic as a small library with R
and Python bindings at 10.5281/zenodo.22168263, so a reader can apply the same two identities to
their own scheme. [No related manuscript is under
consideration elsewhere.]

Yours sincerely,

\vspace{1.5em}

Luciano Moffatt

\end{document}
